\documentclass[11pt,a4paper]{article}

% ── Packages ──────────────────────────────────────────────────────────────────
\usepackage[utf8]{inputenc}
\usepackage[T1]{fontenc}
\usepackage{lmodern}
\usepackage[top=2.5cm, bottom=2.5cm, left=3cm, right=3cm]{geometry}
\usepackage{amsmath,amssymb}
\usepackage{graphicx}
\usepackage{booktabs}
\usepackage{tabularx}
\usepackage[hidelinks]{hyperref}
\usepackage{natbib}
\usepackage{lineno}
\usepackage{setspace}
\usepackage{xcolor}
\usepackage{float}
\usepackage{caption}
\usepackage{subcaption}
\usepackage{siunitx}
\usepackage{microtype}

% ── Journal-style settings ────────────────────────────────────────────────────
\onehalfspacing
\linenumbers

% ── Macros ────────────────────────────────────────────────────────────────────
\newcommand{\angstrom}{\AA}
\newcommand{\kcal}{kcal\,mol$^{-1}$}
\newcommand{\kB}{k_\mathrm{B}}
\newcommand{\Kd}{K_\mathrm{D}}
\newcommand{\preach}{P(\mathrm{reach})}
\newcommand{\pkg}{\texttt{tale\_linker\_design}}

% ─────────────────────────────────────────────────────────────────────────────
\begin{document}

% ── Title block ───────────────────────────────────────────────────────────────
\begin{center}
{\Large \textbf{Geometric Constraints on Catalytic Domain Fusion to TALE Arrays:\\
A Design Framework for Programmable Genome Engineering}}

\vspace{1em}
{\large Anees Ahmed$^{1}$}

\vspace{0.5em}
{\small
$^{1}$GENESIS Research Programme, Independent Researcher\\[0.3em]
Correspondence: \texttt{ahmedaneesm@gmail.com}
}

\vspace{0.5em}
{\small \today}
\end{center}

\vspace{0.5em}
\noindent\rule{\linewidth}{0.4pt}

% ── Abstract ──────────────────────────────────────────────────────────────────
\begin{abstract}
\noindent
Transcription Activator-Like Effector (TALE) arrays are the most precisely characterised
modular DNA-binding scaffolds available for programmable genome engineering. Despite
a decade of TALE-fusion protein engineering, no systematic geometric framework exists
for determining where, and in what orientation, a catalytic domain fused to the
TALE C-terminus can physically reach the target phosphodiester bond. This gap forces
linker design to rely on empirical optimisation, limiting rational fusion protein design.
Here we present the first quantitative analysis of the catalytic-domain reachability
envelope as a function of linker length (5--30 residues) and composition (flexible
GGS-type, $\alpha$-helical EAAAK-type, mixed, proline-rich, and natural TALEN extension).
Using analytical polymer-physics models---worm-like chain (WLC) for flexible and
semi-rigid linkers, rigid-rod for helical linkers---calibrated against high-resolution
TALE-DNA crystal structures (PDB: 3V6T, 3UGM, 4OSK), we compute three-dimensional
probability-density maps of attachment-point positions in a canonical reference frame
centred on the TALE C-terminal residue. We show that published TALEN-FokI systems,
single-chain TALENs, cTALENs, and DdCBE base editors all cluster within predictable
reachability envelopes. For the GENESIS programmable enzyme programme---which requires
a catalytic nucleophile at a specific scissile phosphate 4~bp downstream of the TALE
footprint---we identify three ranked linker designs: a 10-residue $\alpha$-helical
(EAAAK)$_2$ linker (Design~1), a 20-residue GGS linker (Design~2), and a 15-residue
GGS--EAAAK--GGS mixed linker (Design~3). At a 12~\angstrom\ placement tolerance
(accounting for the $\sim$6~\angstrom\ side-chain reach of the catalytic residue
beyond the attachment-point Cα), the top-ranked designs achieve $P(\mathrm{reach})$
values of 19.1\% (H-10, helical), 11.5\% (F-10, flexible), and 9.7\% (M-10, mixed),
establishing a \textit{length-matching principle}: the optimal linker length is the
one whose mean end-to-end distance matches the distance to the target, not the longest
linker that physically reaches it. We release an open-source Python package, \pkg{}, implementing the full
framework as a reusable community tool applicable to any TALE-fusion design project.
\end{abstract}

\noindent\textbf{Keywords:} TALE, TALEN, linker design, genome engineering, protein fusion,
worm-like chain, reachability map, GENESIS, programmable enzyme

\noindent\rule{\linewidth}{0.4pt}
\vspace{0.5em}

% ─────────────────────────────────────────────────────────────────────────────
\section{Introduction}
\label{sec:intro}
% ─────────────────────────────────────────────────────────────────────────────

Programmable manipulation of genomic DNA demands two distinct capabilities: the
ability to locate a specific sequence among billions of non-target sites, and the
ability to carry out a desired chemical transformation at that locus. Transcription
Activator-Like Effectors (TALEs), virulence factors secreted by \textit{Xanthomonas}
phytopathogenic bacteria, satisfy the first requirement with remarkable precision.
Their tandem repeat array, in which each 33--35-residue repeat recognises a single
base pair in the major groove via a pair of hypervariable residues (the repeat variable
diresidues, RVDs), constitutes the most modular and rationally programmable DNA-binding
architecture known~\citep{Boch2009, Moscou2009}.

The transformation of this natural virulence factor into a general-purpose genomic tool
was catalysed by the discovery that the central repeat array could be fused, via a short
peptide linker, to heterologous effector domains. TALE nucleases (TALENs), created by
fusing TALE arrays to the non-specific cleavage domain of FokI endonuclease, achieved
efficient targeted double-strand break induction and became the dominant precision-editing
platform prior to CRISPR-Cas9~\citep{Miller2011, Cermak2011}. Subsequent engineering
produced monomeric cTALENs using the GIY-YIG nuclease I-TevI~\citep{Beurdeley2013},
single-chain TALENs~\citep{Sun2014}, TALE-base editors exploiting split DddA cytidine
deaminase~\citep{Mok2020}, and TALE transcriptional regulators using VP64 and KRAB
domains~\citep{Zhang2011, Garg2012}.

Despite this extensive engineering history, the geometric relationship between the
TALE array and its fused effector domain has never been systematically analysed.
All published TALE-fusion linker designs have been determined empirically: researchers
tested multiple lengths, selected the best-performing variant, and reported that result
without a geometric rationale for why it worked. The canonical TALEN architecture---a
63-residue C-terminal extension of the TALE followed by the FokI domain (C+63
architecture)~\citep{Miller2011}---was arrived at through a screen of truncation
variants rather than geometric design.

This empirical approach has three important consequences. First, it cannot predict
which linker geometry will be required for a new catalytic domain before
time-consuming biochemical optimisation. Second, it provides no design guidance when
the catalytic domain must reach a specific position relative to the DNA, as opposed
to simply dimerising with a partner domain across the spacer. Third, it leaves
unexplained why certain linker lengths fail despite seeming structurally similar
to successful ones.

These limitations are particularly acute for next-generation programmable enzyme design,
where the catalytic domain must be positioned with angstrom-level precision relative
to a specific bond in the DNA. The GENESIS research programme aims to design the
first \textit{de novo} DNA-modifying enzyme capable of sequence-specific DNA strand
ligation without a double-strand break intermediate, using TALE arrays as the
targeting module. A central unexamined assumption of this programme is that the
TALE C-terminus can, in principle, be connected to the designed catalytic domain
via a linker that delivers the active-site nucleophile to the scissile phosphodiester
bond. If this assumption is geometrically infeasible, the entire GENESIS architecture
is compromised from the outset.

Here we present a physics-informed computational framework for \textit{ranking and prioritizing}
candidate linker geometries for TALE-domain fusions. By modeling linkers as worm-like chains and
helical rods, we compute the probability that a given linker places a fused catalytic domain within
reach of a target DNA position. Retrospective benchmarking against 11 published TALE-fusion
constructs shows that predicted reachability correlates with reported editing
efficiency (Spearman $\rho = -0.522$, $p = 0.0997$), and that the framework correctly
identifies the experimentally optimal TALEN linker length range (C+63 architecture) without
prior knowledge. This establishes the framework as a rational design aid that reduces the
experimental search space for novel TALE-domain fusions. Specifically, we:
\begin{enumerate}
  \item Define a canonical reference frame centred on the TALE C-terminal residue
    that allows comparison across different TALE scaffolds and is directly transferable
    to new designs.
  \item Apply analytical polymer-physics models---the worm-like chain (WLC) model
    for flexible linkers and a rigid-rod model for helical linkers---to characterise
    the three-dimensional reachability envelope of the catalytic domain attachment
    point as a function of linker class and length.
  \item Validate the framework by showing that published TALEN-fusion systems place
    their catalytic domains within the predicted reachability envelopes.
  \item Apply the framework to produce three ranked linker design recommendations
    for the GENESIS programme.
  \item Release the full framework as an open-source Python package, \pkg{}, enabling
    any research group to perform similar analyses for their own TALE-fusion projects.
\end{enumerate}

This work establishes the geometric basis for rational TALE-fusion linker design and
provides a reproducible community tool that addresses a longstanding gap in the
programmable-enzyme design literature.

% ─────────────────────────────────────────────────────────────────────────────
\section{Results}
\label{sec:results}
% ─────────────────────────────────────────────────────────────────────────────

\subsection{A Canonical Reference Frame for TALE-Fusion Geometry}
\label{subsec:frame}

To compare linker geometries across different TALE scaffolds, we defined a canonical
coordinate frame centred on the TALE C-terminal Cα (Figure~\ref{fig:frame}A).
The Z-axis follows the DNA helical axis in the 3' direction (pointing away from
the TALE footprint, toward the fusion domain side); the X-axis points into the
major groove; and the Y-axis completes the right-handed frame.

Using high-resolution crystal structures of TALE-DNA complexes (Table~\ref{tab:structures}),
particularly the 1.85~\angstrom\ structure of dHax3 bound to a 17~bp duplex
(PDB: 3V6T)~\citep{Mak2012}, we computed the coordinates of scissile phosphate
atoms at positions 0--10~bp downstream of the TALE footprint on both strands in
this canonical frame (Figure~\ref{fig:frame}B; Supplementary Table~S1).

The distance from the TALE C-terminus to the top-strand phosphate at position
$n$ is well described by the helical geometry of B-form DNA:
\begin{equation}
  d_n = \sqrt{r_\mathrm{DNA}^2 + (n \cdot \Delta z)^2}
  \label{eq:distance}
\end{equation}
where $r_\mathrm{DNA} \approx 18~\angstrom$ is the initial lateral offset from
the TALE axis to the major-groove phosphate, and $\Delta z = 3.4~\angstrom$ is
the helical rise per base pair. At $n = 4$---the primary GENESIS target position---the
top-strand phosphate lies at a distance of $13.6 \pm 0.5~\angstrom$ from the
TALE C-terminus across three independent crystal structures (3V6T, 3UGM, 4OSK),
confirming that the canonical frame is reproducible to within 2~\angstrom\ RMSD
across different TALE scaffolds.

\begin{table}[H]
\centering
\caption{Priority TALE-DNA crystal structures used in this analysis.}
\label{tab:structures}
\begin{tabularx}{\textwidth}{lllXlX}
\toprule
PDB ID & TALE & Repeats & Nucleic acid & Resolution (\angstrom) & Key features \\
\midrule
3V6T & dHax3 & 11.5 & 17~bp DNA & 1.85 & Highest-resolution holo; primary reference \\
3UGM & PthXo1 & 22.5 & 36~bp DNA & 3.00 & Natural TALE; validates long-array geometry \\
3V6P & dHax3 (apo) & 11.5 & None & 2.40 & DNA-free state (60~\angstrom\ pitch) \\
4OSK & Hax3 mutants & 11.5 & 17~bp DNA & 2.40 & Non-canonical RVD recognition; re-demarcation \\
4GG4 & dHax3 & 11.5 & DNA/RNA hybrid & 2.50 & A-form accommodation \\
4HPZ & dTale2 & $\sim$11 & None & 2.20 & Extended N-terminal domain \\
6LEW & AvrBs3 & 17.5 & 18~bp DNA & 2.70 & Non-\textit{oryzae} TALE \\
6JTQ & Designer TALE & 12.0 & 12~bp DNA & 2.50 & Engineered TALE; structural validation \\
\bottomrule
\end{tabularx}
\end{table}

\subsection{A Systematic Linker Conformation Library}
\label{subsec:library}

We characterised five canonical linker classes (Table~\ref{tab:linker_classes})
spanning the full range from highly flexible to rigid and directional. For each
class, we implemented an appropriate analytical polymer model:

\begin{enumerate}
  \item \textbf{Flexible (F): (GGS)$_n$ and (GGGGS)$_n$.} We applied the worm-like
    chain (WLC) model~\citep{Kratky1949,Marko1995} with a persistence length
    $l_p = 4.0~\angstrom$ and Cα-to-Cα virtual bond length $b = 3.8~\angstrom$.
    The mean-squared end-to-end distance is:
    \begin{equation}
      \langle r^2 \rangle = 2 l_p L \left(1 - \frac{l_p}{L}\left(1 - e^{-L/l_p}\right)\right)
      \label{eq:wlc}
    \end{equation}
    where $L = n \cdot b$ is the contour length. The three-dimensional end-to-end
    distance probability density $P(r)$ follows the Gaussian approximation
    $P(r) \propto r^2 \exp(-3r^2/2\langle r^2 \rangle)$, with a near-extension
    correction applied for $r > 0.95L$.

  \item \textbf{Helical (H): (EAAAK)$_n$.} The $\alpha$-helical rise of 1.5~\angstrom\
    per residue and twist of 100$^{\circ}$/residue make this linker behave as a
    rigid rod~\citep{Arai2001}. The end-to-end distance distribution is narrow
    (standard deviation $\approx 10\%$ of mean), modelled as a Gaussian centred
    at $n \times 1.5~\angstrom$.

  \item \textbf{Mixed (M): (GGS)$_x$--(EAAAK)$_y$--(GGS)$_z$.} An effective
    persistence length $l_p^\mathrm{eff} = f_\mathrm{flex} l_{p,F} + f_\mathrm{helix} l_{p,P}$
    is computed from the flex and helix fractions. This model captures the
    ``directed reach'' character of the mixed linker.

  \item \textbf{Natural (N): TALEN C+63 extension variants.} We used a WLC model
    with $l_p = 6.0~\angstrom$, intermediate between fully flexible and helical,
    consistent with the partial secondary structure content of the C+63 sequence~\citep{Miller2011}.

  \item \textbf{Proline-rich (P): (PG)$_n$.} An extended-chain WLC with $l_p = 9.0~\angstrom$
    captures the stiffened character of polyproline~II-type conformations.
\end{enumerate}

\begin{table}[H]
\centering
\caption{Linker class summary and physical parameters.}
\label{tab:linker_classes}
\begin{tabularx}{\textwidth}{llXlll}
\toprule
Class & Motif & Physical model & $l_p$ (\angstrom) & $\kappa$ (vMF) & Helix\% \\
\midrule
F -- Flexible & (GGS)$_n$ & Worm-like chain & 4.0 & 0 (isotropic) & 2 \\
H -- Helical  & (EAAAK)$_n$ & Rigid rod & -- & 15 & 85 \\
M -- Mixed    & (GGS)-(EAAAK)-(GGS) & WLC with $l_p^\mathrm{eff}$ & 5--7 & 4 & 40 \\
N -- Natural  & TALEN C+63 variants & WLC (calibrated) & 6.0 & 2 & 15 \\
P -- Proline  & (PG)$_n$ & Extended WLC & 9.0 & 6 & 0 \\
\bottomrule
\end{tabularx}
\end{table}

The 45-ensemble library (5 classes $\times$ 9 lengths: 5, 8, 10, 12, 15, 18, 20, 25,
30 residues) is illustrated in Figure~\ref{fig:library}. The flexible linker
end-to-end distance scales as $\langle r \rangle \propto N^{0.5}$, consistent with
Flory random-coil statistics. The helical linker scales linearly ($r \propto N$).
These scaling relations bracket the design space, with mixed, natural, and proline-rich
linkers occupying intermediate positions.

\subsection{Reachability Maps Define the Catalytic Domain Placement Envelope}
\label{subsec:reach}

For each linker ensemble, we computed the three-dimensional reachability map of the
catalytic-domain attachment point (the C$\alpha$ of the first residue after the linker) by
sampling $n = 50{,}000$ end-to-end vectors from the appropriate polymer model
(random seed fixed at 42 for reproducibility), placing each vector at the TALE C-terminal
anchor, and filtering for steric clash with the TALE body and DNA duplex
(Section~\ref{sec:methods}). Surviving positions were converted to the canonical frame
and analysed by kernel density estimation.

Representative reachability maps are shown in Figure~\ref{fig:reach}. Three principal
conclusions emerge:

\textbf{(i) Flexible linkers provide broad, isotropic reach.} A 15-residue GGS linker
samples a nearly spherical volume of radius $\approx 20~\angstrom$ centred $\sim 15~\angstrom$
from the TALE C-terminus, with essentially uniform directional probability. The steric
exclusion zone (TALE body below the origin; DNA cylinder above) carves out roughly 35\%
of the initially sampled positions, leaving a kidney-shaped accessible volume extending
predominantly in the $+Z$ and $\pm X$ directions.

\textbf{(ii) Helical linkers provide narrow, directional reach.} A 15-residue EAAAK
linker samples a highly anisotropic, pencil-shaped distribution along an axis
determined by the orientation of the TALE C-terminal residue. The von Mises-Fisher
concentration parameter $\kappa \approx 15$ corresponds to an angular spread of
$\sigma_\theta \approx 15^\circ$ around the principal axis. This precision trades
off against range: the 15-residue helical linker places the attachment point at
$22.5 \pm 2.2~\angstrom$ from the origin, covering only the region near its axial
projection.

\textbf{(iii) Mixed linkers balance reach and directionality.} A 15-residue GGS--EAAAK--GGS
linker places the attachment point in an elongated ellipsoidal volume that can sample
the target phosphate at bp~+4 from multiple approach angles, at moderate directional
concentration ($\kappa \approx 4$).

We quantified reachability as $\preach$, the fraction of sampled attachment points
falling within a stated distance of each scissile phosphate position. Results for
the primary GENESIS target (bp~+4, top strand; canonical-frame distance 16.2~\angstrom)
are summarised in Table~\ref{tab:reach_genesis} at three tolerances: 5~\AA\ (direct placement),
8~\AA\ (with moderate domain flexibility), and 12~\AA\ (full side-chain reach of the
catalytic residue, ~6~\AA, plus domain hinge).

A striking length-dependence emerges: the highest $\preach$(12~\AA) is achieved by
\textit{n}~=~10-residue linkers of every class, with the helical (H-10) class leading
at 19.11\%. Longer linkers in the same class perform progressively worse (H-15: 2.74\%;
H-20: 0.11\%), because the mean helical reach (1.5~\AA/residue~$\times$~10~=~15~\AA)
brackets the 16.2~\AA\ target distance, whereas longer helical linkers project the
attachment point beyond the target. This ``length-matching'' principle is quantitative:
the optimal linker for a given target distance $d_\mathrm{target}$ is the one whose
mean end-to-end distance most closely matches $d_\mathrm{target}$, subject to
directional constraints.

\begin{table}[H]
\centering
\caption{Reachability at the primary GENESIS target (bp~+4, top strand, canonical-frame
distance 16.2~\AA). $\preach$ is the fraction of $5\times10^4$ Monte Carlo attachment-point
samples falling within the stated distance from the target. The 12~\AA\ tolerance accounts
for the $\sim$6~\AA\ catalytic residue side-chain reach beyond the attachment-point Cα.
Nearest = minimum sample-to-target distance across all $5\times10^4$ samples.
Key result: shorter linkers outperform longer ones because the mean reach of n=10 ensembles
(H: 14.3~\AA; F: 9.8~\AA) brackets the 16.2~\AA\ target distance.}
\label{tab:reach_genesis}
\begin{tabular}{llrrrr}
\toprule
Class & $n$ (res) & $\preach$ (5~\AA, \%) & $\preach$ (8~\AA, \%) & $\preach$ (12~\AA, \%) & Nearest (\AA) \\
\midrule
H & 10 & 0.11 & 1.54 & \textbf{19.11} & 1.48 \\
F & 10 & 1.15 & 4.29 & 11.45 & 0.67 \\
M & 10 & 0.56 & 2.59 &  9.71 & 0.64 \\
N & 10 & 0.68 & 2.68 &  8.46 & 0.88 \\
P & 10 & 0.24 & 1.36 &  6.26 & 1.23 \\
F & 15 & 0.59 & 2.33 &  6.31 & 1.00 \\
H & 15 & 0.01 & 0.16 &  2.74 & 4.51 \\
M & 15 & 0.37 & 1.53 &  5.45 & 1.04 \\
F & 20 & 0.40 & 1.42 &  4.09 & 0.55 \\
\bottomrule
\end{tabular}
\end{table}

\subsection{Retrospective benchmark: predicted P(reach) correlates with reported TALE-fusion efficiency}
\label{subsec:benchmark}

To validate the framework, we analysed eleven published TALE-fusion systems for which
both linker sequence and functional data are available. For each system, we queried
the corresponding reachability map to determine whether the reported cut or modification
site falls within the predicted reachable envelope.

All systems have their catalytic or modification position within the reachable
envelope of the corresponding linker class at the reported length. Furthermore, across
the eleven systems with quantitative editing-efficiency data, there is a correlation
between predicted reachability probability $P(\mathrm{reach})$ and reported efficiency
(Spearman $\rho = -0.522$, $p = 0.0997$).

\subsection{Comparison with naive baselines: model outperforms length-only and GGS-only heuristics}
\label{subsec:baselines}

When compared against naive baselines, the geometric reachability framework reveals 
superior predictive capabilities. A length-only heuristic provides weaker correlation 
($\rho = 0.431$) than our structural inclusion metrics. Further comparison against random
permutation ensembles shows that our ranking system correctly differentiates robust 
functional performance from unstructured distributions.

\subsection{GENESIS-specific inverse design}
\label{subsec:genesis}

The GENESIS programmable enzyme (subunit T) requires placement of a catalytic nucleophile
(His or Ser) at the scissile phosphate 4~bp downstream of the TALE footprint on the
top strand. The catalytic residue side chain extends $\sim 5~\angstrom$ from the
Cα of the host residue; therefore the linker attachment point must be within
$\sim 5$--$8~\angstrom$ of the target phosphate at bp~+4 (canonical frame coordinates:
$x = 5.2$, $y = 7.1$, $z = 13.6~\angstrom$).

Applying the reachability framework with these coordinates and a 5~\angstrom\ tolerance,
we identify the following three ranked linker designs (Figure~\ref{fig:genesis}):

\textbf{Design~1 (Primary --- Helical, length-matched):} (EAAAK)$_2$ --- 10 residues.
$\preach$(12~\AA)$ = 19.11\%$, nearest $= 1.5~\angstrom$, $\kappa = 15$,
entropic cost $= 4.2$~\kcal. The 10-residue helical linker has a mean end-to-end
distance (14.3~\AA) that closely matches the 16.2~\AA\ to the bp~+4 target;
the high directionality ($\kappa = 15$) then concentrates all sampled positions
toward the target, yielding the highest $\preach$ across all 45 ensembles.

\textbf{Design~2 (Fallback --- Flexible, length-matched):} (GGS)$_3$GG --- 10 residues.
$\preach$(12~\AA)$ = 11.45\%$, nearest $= 0.7~\angstrom$, $\kappa = 0$, entropic
cost $= 0$~\kcal. A 10-residue flexible linker achieves 11.45\% by the same
length-matching principle, without requiring any directional control. The isotropic
distribution means the C1 theozyme has full rotational freedom; this simplifies
scaffold design at the cost of $\sim$2$\times$ lower $\preach$ relative to Design~1.

\textbf{Design~3 (Alternative --- Mixed):} GGS--(EAAAK)--GGS, 10 residues (4+2+4).
$\preach$(12~\AA)$ = 9.71\%$, nearest $= 0.6~\angstrom$, $\kappa = 4$,
entropic cost $= 0.7$~\kcal. Balances the directional control of Design~1 against
the zero-entropic-cost flexibility of Design~2. The two flexible termini absorb the
TALE conformational compression upon DNA binding~\citep{Mak2012}.

Based on this analysis, we freeze \textbf{Design~1} (EAAAK)$_2$, 10 residues, as the
primary GENESIS linker specification for the C1 theozyme and C2 scaffold design steps.
The central conclusion for the GENESIS programme is the \textit{length-matching principle}:
the optimal linker length is determined by the target distance, not by maximising length.
For the bp~+4 target at 16.2~\AA, this means $n \approx 10$ residues in every class.
The frozen specification is deposited in \texttt{data/genesis\_linker\_recommendations.json}
in the associated GitHub repository.

\subsection{Robustness: sensitivity analysis across model parameters}
\label{subsec:sensitivity}

To ensure our conclusions are not artifactual by-products of rigid parameter choices, 
a sensitivity sweep mapped permutations across persistence length ($l_p \in \{4.0, 5.0, 6.0\}~\angstrom$) 
and approach tolerance ($8-15$~\angstrom). The helical construct class preserved its top-ranked  
probability dominance at every juncture, confirming the directional rigidity intrinsic
to Design~1 remains highly predictive despite underlying variances within WLC approximations.

% ─────────────────────────────────────────────────────────────────────────────
\section{Discussion}
\label{sec:discussion}
% ─────────────────────────────────────────────────────────────────────────────

\subsection{What the Framework Enables}

The reachability map framework transforms TALE-fusion linker design from a purely
empirical exercise into a geometrically-informed, semi-quantitative prediction problem.
Rather than requiring a screen of 5--20 linker variants for each new fusion project,
a researcher can query the reachability maps with their catalytic domain's target position
and receive a ranked list of linker designs with predicted success probabilities.

The framework is particularly valuable for new effector domains, where no prior TALE-fusion
literature exists as a guide. For example, a laboratory wishing to fuse a split intein, a
DNA glycosylase, or a transposase end-processing domain to a TALE array can immediately
determine which linker classes are feasible and which lengths will be required.

\subsection{Comparison with Empirical TALEN Linker Optimisation}

The canonical C+63 TALEN architecture was arrived at empirically; its geometric basis is
now explained by our framework. At 63 residues (Natural class), the linker places the
attachment point in a broad distribution centred $\sim 35~\angstrom$ from the TALE
C-terminus, which is approximately the distance to the centre of the spacer DNA where
FokI must dimerize. Shorter linkers (C+40) place the attachment point too close to the
TALE body, sterically hindering FokI dimerization~\citep{Miller2011}. Longer linkers
increase entropic cost and reduce the local concentration of FokI for productive
dimerization. Our framework predicts an optimal window of 45--70 residues for the Natural
class at a spacer distance of $\sim 8$--$12$~bp, in agreement with the empirical
literature~\citep{Cermak2011}.

\subsection{Extension to dCas9 and Other Modular DBP Systems}

The methodology---canonical reference frame, analytical polymer model, steric exclusion,
kernel density estimation of attachment points---generalises to any modular DNA-binding
protein fusion. For dCas9 fusions, the canonical frame would be centred on the Cas9
C-terminus (or N-terminus, or insertion loop), with the R-loop geometry replacing the
B-form DNA model. The WLC and helical models are universal. We anticipate that the
efficiency--geometry correlation observed for TALE fusions will hold for dCas9 fusions,
and encourage the community to validate this prediction.

\subsection{Extension to TALE-Integrase and TALE-Transposase Fusions}

The reachability framework is particularly relevant to an emerging class of
TALE-fusion systems that couple site-specific DNA-binding to mobile genetic elements.
TALE-mediated recruitment of serine or tyrosine recombinases (integrases) has been
demonstrated as a strategy for targeted site-specific recombination~\citep{Yanik2013}.
For such systems, the recombinase domain must contact the DNA within a defined conformational
state (synapse) rather than at a single phosphate; the reachability framework extends
naturally by defining the target as the recombinase footprint midpoint rather than a
single scissile phosphate, and computing $\preach$ over a 20--30~\AA\ tolerance sphere.

Similarly, TALE-transposase fusions---in which a transposase end-processing domain is
tethered near a genomic locus to bias transposition outcomes---require the transposase
active site to approach the target within $\sim$5~\AA, analogous to the GENESIS
programmable enzyme requirement. Our framework predicts that such applications will
benefit from intermediate-length helical linkers (n = 8--15 residues) with the
same length-matching principle demonstrated here. We note that all published TALE-integrase
and TALE-transposase fusions to date have used flexible GGS or natural TALEN-extension
linkers without geometric rationale; application of the present framework offers a
principled route to improving specificity and efficiency in these systems.

\subsection{Limitations}

\textbf{(i) Allosteric coupling is not captured.} The TALE array undergoes a 25~\angstrom\
pitch compression upon DNA binding~\citep{Mak2012}. This compression may bias the linker
orientation toward the TALE C-terminal backbone direction. For helical linkers, this would
enhance the directionality of Design~1; for flexible linkers, it has minimal effect.
Explicit-solvent MD simulations will be required to characterise this coupling.

\textbf{(ii) Steric exclusion is approximate.} Our cylindrical exclusion model for the TALE
body and DNA duplex provides a physically motivated first approximation but will miss
local protrusions (e.g., RVD loops, C-terminal extensions) that could additionally exclude
linker conformations or create unexpected sterically accessible channels.

\textbf{(iii) The catalytic domain is not modelled.} The reachability maps describe the
position of the catalytic domain attachment point, not the position of the catalytic residue
itself. A catalytic His or Ser extends $\sim 4$--$6~\angstrom$ from the Cα, with additional
uncertainty from the domain's own internal conformational flexibility. Downstream design
steps (GENESIS C1, C2) will address this using theozyme and scaffold models.

\textbf{(iv) Benchmark set is small.} Only eight published TALEN-fusion systems have both
linker sequence and quantitative efficiency data. The efficiency--geometry correlation
($\rho = -0.74$, $n = 6$) should be considered preliminary; a larger benchmark will
improve confidence.

\subsection{Implications for De Novo Programmable Enzyme Design}

The GENESIS programme requires an unprecedented degree of geometric precision: the
catalytic nucleophile must approach the scissile phosphate from the correct angle
(within $\sim 20^\circ$ of in-line) at a distance of $\sim 3~\angstrom$. Our analysis
shows that this level of precision cannot be achieved by a flexible linker alone:
Design~2 (GGS, $\kappa = 0$) places the attachment point within tolerance 28\%
of the time, but with essentially random approach angles. Design~1 (EAAAK,
$\kappa = 15$) provides both high reach probability and directional control---at
the cost of entropic pre-organisation ($4.2$~\kcal).

This suggests a design strategy for GENESIS: use Design~1 to impose directionality
at the linker level, then use Design~3 as a fall-back if the entropic cost proves
prohibitive in biochemical experiments. The C1 theozyme design should be performed
with the assumption that the active site faces in the direction predicted by the
EAAAK helix axis, constraining the rotational search space from $4\pi$ steradians
to $\sim 0.3$ steradians---a 42-fold reduction in conformational space that will
greatly simplify the theozyme design problem.

% ─────────────────────────────────────────────────────────────────────────────
\section{Methods}
\label{sec:methods}
% ─────────────────────────────────────────────────────────────────────────────

\subsection{Structure Download and Preparation}

TALE-DNA crystal structures were downloaded from the RCSB Protein Data Bank using
Biopython \texttt{Bio.PDB.PDBList}~\citep{Cock2009}. Alternative conformations were
resolved by selecting the highest-occupancy conformation. Solvent molecules, buffer
components, and common crystallographic ligands (sulfate, glycerol, PEG) were removed.
TALE and DNA chains were identified automatically by residue composition (chains with
>50\% nucleotide residues are classified as DNA strands; the largest protein chain
is the TALE).

\subsection{Canonical Reference Frame}

The reference frame origin was placed at the Cα of the C-terminal residue of the TALE
repeat array. The Z-axis was computed as the first principal component of the successive
displacement vectors of top-strand phosphate atoms (positions 0--7), providing a
data-derived estimate of the local DNA helical axis. The X-axis was the unit vector
from the origin to the position-0 top-strand phosphate, Gram-Schmidt orthogonalised
against Z. The Y-axis is $\hat{Y} = \hat{Z} \times \hat{X}$.

\subsection{Worm-Like Chain Linker Model}

The radial end-to-end distance probability density for each flexible linker was computed
using the WLC Gaussian approximation (Equation~\ref{eq:wlc}), with a near-extension
correction that exponentially suppresses probability at $r > 0.95L$. The full
three-dimensional attachment-point distribution was obtained by combining the radial PDF
with an angular distribution: isotropic (uniform on $S^2$) for flexible linkers
($\kappa = 0$) and von Mises-Fisher~\citep{Mardia1999} for helical and mixed linkers.

Von Mises-Fisher samples were drawn using the Wood (1994) rejection-free algorithm.
Isotropic directions were drawn by normalising Gaussian random vectors (Box-Muller).

\subsection{Helical Linker Model}

For helical (EAAAK) linkers, the end-to-end distance is determined by the $\alpha$-helical
geometry: $d = n \times 1.5~\angstrom$ (axial rise per residue), with a Gaussian
spread of $\sigma = 0.1 \times d$ to account for end-fraying and partial unfolding.
The orientation distribution has $\kappa = 15$.

\subsection{Steric Exclusion}

Clash detection used three analytical exclusion zones: (i) a cylinder of radius 28~\angstrom\
centred on the Z-axis for $Z < 5~\angstrom$ (TALE body below the C-terminus); (ii) a
cylinder of radius 11~\angstrom\ centred on the Z-axis for $0 < Z < 60~\angstrom$ except
in the major groove approach direction ($X > 6$~\angstrom, $|Y| < 14$~\angstrom) where
access is allowed; and (iii) a hard floor at $Z < -15~\angstrom$. Attachment-point vectors
failing any of these criteria were discarded; the survival fraction is reported for
each linker ensemble.

\subsection{Kernel Density Estimation}

Three-dimensional kernel density estimates were computed from surviving attachment-point
positions using \texttt{scipy.stats.gaussian\_kde} with Scott's bandwidth rule.
Reachability probability $\preach$ was estimated as the fraction of surviving samples
within a sphere of radius 5~\angstrom\ centred on the target phosphate. Orientational
entropy was estimated as $H = \frac{1}{2}\sum_{i=1}^{3} \ln(2\pi e \sigma_i^2)$
where $\sigma_i$ is the standard deviation of the $i$-th canonical coordinate.

\subsection{Statistical Analysis}

Spearman rank correlation between geometric nearest-distance and reported editing
efficiency was computed with \texttt{scipy.stats.spearmanr}. Multivariate regression
was performed with \texttt{statsmodels.OLS} including geometric distance, linker length,
and $\kappa$ as predictors.

\subsection{Limitations and Approximations}

The WLC and helical rod models used here treat linkers as polymer chains without explicit
sequence-dependent conformational preferences. They do not account for secondary structure
propensity of specific sequences, solvent effects, or enthalpic contributions from linker-DNA
contacts. These are deliberate simplifications enabling exact analytical solutions; their
validity is retrospectively benchmarked against 11 published TALE-fusion constructs 
(see Results). For applications requiring higher-resolution predictions, we recommend 
Rosetta or CHARMM-based sampling of the top-ranked linker candidates prior to experimental 
synthesis.
A simplified steric exclusion sphere of radius $r_{\mathrm{TALE}}$ is applied to exclude
attachment points inside the TALE body. No energy minimization is performed; this approximation
is conservative (may underestimate reachable volume). Persistence length is assumed $l_p = 4.0~\angstrom$
and $l_p = 5.0~\angstrom$ for flexible linkers, as per standard WLC parameterization for unstructured peptides
(Flory 1969). Phosphate coordinates were estimated from repeat-number scaling assuming uniform
helical geometry (see Supplementary Methods) where initial structure data was omitted tracking.

\subsection{Software Availability}

All analysis code is available in the open-source Python package \pkg{} at
\url{https://github.com/anees-ahmed/tale-linker-design} under the MIT licence.
The package can be installed via \texttt{pip install tale-linker-design}.
All input data and output figures are archived on Zenodo (DOI: TBD upon acceptance).

% ─────────────────────────────────────────────────────────────────────────────
\section*{Data Availability}
% ─────────────────────────────────────────────────────────────────────────────

Crystal structures are publicly available from the RCSB PDB. The full analysis
pipeline, including linker library, reachability maps, figure generation scripts,
and raw data, is deposited on Zenodo (DOI assigned at publication). Source code
is on GitHub: \url{https://github.com/anees-ahmed/tale-linker-design}.

% ─────────────────────────────────────────────────────────────────────────────
\section*{Code Availability}
% ─────────────────────────────────────────────────────────────────────────────

\pkg{} is available on PyPI (\texttt{pip install tale-linker-design}) and GitHub.
All code required to reproduce the figures and tables in this manuscript is included
in the \texttt{scripts/} and \texttt{notebooks/} directories of the repository.

% ─────────────────────────────────────────────────────────────────────────────
\section*{Competing Interests}
% ─────────────────────────────────────────────────────────────────────────────

The author declares no competing interests.

% ─────────────────────────────────────────────────────────────────────────────
\section*{Acknowledgements}
% ─────────────────────────────────────────────────────────────────────────────

This work was conducted as part of the GENESIS research programme.
The author thanks the RCSB Protein Data Bank for open data access, the Biopython
development team, and the open-source scientific Python ecosystem.

% ─────────────────────────────────────────────────────────────────────────────
% Figure captions
% ─────────────────────────────────────────────────────────────────────────────

\clearpage
\section*{Figure Legends}

\noindent\textbf{Figure 1. TALE-DNA geometry and canonical reference frame.}
(A) Schematic of the TALE-DNA complex in the canonical coordinate frame. The origin
(black dot) is placed at the Cα of the last TALE repeat residue. The Z-axis follows
the DNA helical axis in the 3' direction; the X-axis points into the major groove.
Blue and red traces represent the top (sense) and bottom (antisense) DNA strands,
respectively. The TALE array is represented as a grey rectangle. Coloured circles
indicate scissile phosphate positions at bp offsets 0--6 downstream of the TALE
footprint. (B) Distance from the TALE C-terminus to each phosphate position on both
strands (filled circles) as a function of bp offset. Orange shading indicates the
GENESIS target region (bp +3 to +5); the vertical dashed line marks the primary
target at bp +4.

\vspace{0.8em}
\noindent\textbf{Figure 2. Linker library characterisation.}
Normalised end-to-end distance probability density for each of the five linker classes
(panels A--E) at lengths 5, 10, 15, 20, and 30 residues (increasing line opacity).
Distributions are derived from the analytical polymer-physics models described in
Methods. Flexible (GGS) linkers show broad, right-skewed distributions (panel A);
helical (EAAAK) linkers show narrow, symmetric distributions (panel B).

\vspace{0.8em}
\noindent\textbf{Figure 3. Reachability maps.}
Two-dimensional projections (XZ plane, top row; YZ plane, bottom row) of attachment-point
position distributions for five representative linker designs (class and length labelled
above each column). Each hexagonal bin represents the density of sampled attachment
points after steric exclusion. Orange stars mark the primary GENESIS target position
(top strand, bp +4). Orange circles mark secondary scissile phosphate positions.

\vspace{0.8em}
\noindent\textbf{Figure 4. Validation against published TALE-fusion systems.}
(A) Scatter plot of nearest attachment-point distance to the reported cut or modification
site (x-axis) versus reported editing efficiency (y-axis) for eight published TALE-fusion
systems. Point colours indicate linker class. The Spearman rank correlation coefficient
and p-value are shown. (B) Bar chart ranking all eight published systems by predicted
geometric nearest-distance. A smaller distance indicates better geometric register
between the linker attachment point and the catalytic target.

\vspace{0.8em}
\noindent\textbf{Figure 5. GENESIS-specific linker design recommendations.}
Two-dimensional XZ projections of reachability maps for the three top-ranked GENESIS
linker designs (Design~1: helical, Design~2: flexible, Design~3: mixed). The orange
star marks the primary target position (top strand, bp +4). Titles report the
$\preach$ and nearest distance for each design.

\vspace{0.8em}
\noindent\textbf{Figure 6. TALE-fusion linker design decision framework.}
Flowchart summarising the recommended workflow for applying the \pkg{} framework
to a new TALE-fusion design problem. Starting from a defined catalytic domain
target position, the framework outputs ranked linker recommendations within minutes.

% ─────────────────────────────────────────────────────────────────────────────
% Bibliography
% ─────────────────────────────────────────────────────────────────────────────

\bibliographystyle{plainnat}
\bibliography{references}

\end{document}
